\documentclass[12pt]{article}

\usepackage[paper=letterpaper,margin=2.5cm]{geometry} % Set Margins

%% Math and math fonts
\usepackage{amsmath, amsthm, amssymb, amsfonts}
\usepackage{bbold} % for \mathbbm{1}

% date
\usepackage[nodayofweek]{datetime}

% Color
\usepackage{color, xcolor}

% Misc
\usepackage{environ}  % \collect@body in asmmath
\usepackage{graphicx} % \includegraphics options
\usepackage{mdframed} % text boxes
\usepackage{indentfirst} % Indent first paragraph after section header
\usepackage[shortlabels]{enumitem} % Control enumerate items with [(a)]
\usepackage{comment} % Comments
\usepackage{fancyhdr} % Headers and footers

% Tables
\usepackage{array}

% Sub-figures and figure placement
\usepackage{caption}
\usepackage{subcaption}
\usepackage{float} 

% Graphing
\usepackage{pgfplots}
\pgfplotsset{compat=1.17}
\usepackage{tikz}

% Title Placement
\usepackage{titling}
\setlength{\droptitle}{-6em}

%set indent to 
\setlength{\parindent}{0pt}
\setlength{\parskip}{7pt}


% % Move the header up
\usepackage{titling}
\posttitle{\par\centering}
\preauthor{\centering\normalsize}
\postauthor{\par\centering\vspace{-1em}}
\predate{\centering\normalsize}
\postdate{\par\centering\vspace{-1em}}

% ----------------------------------------
% TITLE
% ----------------------------------------

\title{Rules of Derivatives}
\date{}

\begin{document}
\maketitle

\textbf{1. Constant Rule:} The derivative of a constant $k$ is zero, \(\frac{d}{dx}k = 0
\).\\

\textbf{2. Power Rule:} The derivative of $kx^n$, where $k$ and $n$ are constants, is \(
\frac{d}{dx}(kx^n) = knx^{n-1}
\).\\

\textbf{3. Sum Rule:} The derivative of a sum is the sum
of the derivatives, \(
\frac{d}{dx}[f(x)+g(x)]
= f'(x)+g'(x)
\).\\

\textbf{4. Product Rule:} The derivative of a product is
\(
\frac{d}{dx}[f(x)g(x)]
= f'(x)g(x)+f(x)g'(x)
\).\\

\textbf{5. Quotient Rule:} The derivative of a quotient is \(
\frac{d}{dx}
\left[\frac{f(x)}{g(x)}\right]
=
\frac{f'(x)g(x)-f(x)g'(x)}
{[g(x)]^2} \text{, } g(x) \neq 0
\).\\

\textbf{6. Chain Rule:} The derivative of a function nested inside another function is
\[
\frac{d}{dx}[f(g(x))]
=
f'(g(x))g'(x).
\]

\textbf{7. Natural Log Rule:} The derivative of the natural log of a function is
\[
\frac{d}{dx}\ln(f(x))
=
\frac{f'(x)}{f(x)} \quad \text{where } f(x) >0.
\]

\textbf{8. Exponent Rule:} The derivative of an exponential function with base $e$ is
\[
\frac{d}{dx}e^{f(x)}
=
f'(x)e^{f(x)}.
\]

\textbf{9. Total Derivative: } For a function of two variables, $F(x,y)$,
the total derivative is
\[
dF(x,y)
=
F_x\,dx+F_y\,dy \quad \text{or equivalently,} \quad 
dF(x,y)
=
\frac{\partial F}{\partial x}\,dx
+
\frac{\partial F}{\partial y}\,dy.
\]

The symbol $\partial$ denotes a
partial derivative. Partial derivatives measure how a function
changes with respect to one variable while holding the other variables constant.\\

\textbf{10. Implicit Function Theorem:} 

Suppose $F(x,y)$ is held constant along a level set, so $F(x,y)=c$. Because $F$ is constant along this curve,
its total derivative is 0.

\[
dF(x,y)=F_x\,dx+F_y\,dy=0.
\]

With some algebra,
\[
F_y\,dy=-F_x\,dx \implies 
\frac{dy}{dx}
=
-\frac{F_x}{F_y} \quad \text{or equivalently,} \quad 
\frac{dy}{dx}
=
-\frac{\frac{\partial F}{\partial x}}
{\frac{\partial F}{\partial y}}
\]


\end{document}
